\documentclass[12pt]{article}
\usepackage{graphicx, geometry, hyperref} % Required for inserting images
\usepackage{physics,caption, amsmath, amsfonts, tikz, minted}
\usepackage{scalerel, cases, multirow, multicol}
\usepackage{tabularx}
\graphicspath{{images/}}
\hypersetup{colorlinks=true,citecolor=black,linkcolor=black,urlcolor=black}
\geometry{
a4paper,
total={170mm,257mm},
left=20mm,
top=20mm,
}
% \usepackage[table]{xcolor}
% \usepackage{eulervm}
% \usepackage{booktabs}


%-------------new commands----------------------
\newcommand{\bc}[1]{{\color{blue}{#1}}}
\newcommand{\red}[1]{{\color{red}{#1}}}
\begin{document}

\begin{center}
    \Large \bfseries{Weekly Report}
\end{center}

\begin{center}
    \Large \bfseries{CS3500: Operating Systems}
\end{center}


\begin{center}
    \Large \bfseries{Visualisation Tool for Process Scheduling}
\end{center}
\begin{center}
    \vspace{1.2cm}
    \includegraphics[width=5cm]{logos and images/IIT_Madras_Logo.svg.png}

    \textbf{Computer Science and Engineering}
    \vspace{0.3cm}

    \textbf{Indian Institute of Technology  Madras}
    \vspace{0.3cm}

    \textbf{Jul - Nov 2024}
    \vspace{1.2cm}

    \noindent \textbf{Under the supervision}
    \vspace{0.2cm}

    \noindent \textbf{of}
    \vspace{0.2cm}

    \noindent\textbf{Prof. Janakiram D}
     \vspace{0.2cm}

    \noindent \textbf{submitted by}
    \vspace{0.3cm}

    \noindent \textbf{Team 8}
    \begin{multicols}{2}
        \begin{itemize}
            \item Anjali Samudrala (CS22B046)
            \item Chaitanya Sai Teja G (CS22B036)
            \item Jwala Likitha Reddy M (CS22B078)
            \item Karthikeya P (CS22B026)
            \item Naveen Koushik Reddy E (CS22B006)
            \item Navya Sree B (CS22B045)
            \item Rushi Babu G (CS22B040)
            \item Sravya Rangu (CS22B044)
            \item Yashwanth Sai P (CS22B002)
            \item Yaswanth Sai V (CS22B043)
        \end{itemize}
    \end{multicols}
    \vspace{0.5cm}

\end{center}
\clearpage

\section{Frontend Team}
\subsection{Tasks}
\begin{itemize}
    \item The tasks for this week is to identify the components that are required for the frontend. This includes the following.
    \begin{itemize}
        \item Coming up with a frontend design and layout.
        \item Ideation on what all different plots can be displayed.
        \item Informing the backend team about the requirements of various informations.
    \end{itemize}
\end{itemize}

\subsection{Progress}
\subsubsection{Frontend Design}
\section{Backend Team}
\subsection{Linking Backend and Frontend}
The objective is to create a backend server that fetches real-time system process statistics using the \texttt{pidstat} command and streams the data to a frontend via WebSockets. The application uses Flask as the web framework and Flask-SocketIO to establish real-time communication between the backend and frontend.

\subsubsection{Tools and Libraries Used}

\begin{itemize}
    \item \textbf{Flask}: A lightweight web framework for Python that simplifies the development of web applications.
    \item \textbf{Flask-SocketIO}: An extension for Flask that enables real-time communication between the server and the client using WebSockets.
    \item \textbf{subprocess}: A Python module used to run external commands. In this case, it is used to execute the \texttt{pidstat} command, which collects CPU statistics for processes running on the system.
    \item \textbf{re (Regular Expressions)}: A Python module for matching patterns in strings. It is used to parse the output of the \texttt{pidstat} command.
    \item \textbf{threading}: A Python module used to create background threads. In this project, it allows the data-fetching process to run concurrently with the main server.
\end{itemize}

\subsubsection{Architecture}

\begin{itemize}
\item {\bf Backend (Flask Server with SocketIO) :}
The backend is responsible for fetching real-time process statistics using the \texttt{pidstat} command, parsing the output, and sending the data to the frontend using WebSockets.

\item{\bf Real-Time Data Fetching :}
\begin{itemize}
    \item The \texttt{pidstat} command is used to gather CPU statistics for processes every second. The output contains several fields, including process ID (PID), user and system CPU usage, and the process command.
    \item A background thread is created to run the \texttt{pidstat} command continuously, fetching data at specified intervals (e.g., every 2 seconds).
    \item The data is parsed using a regular expression, and relevant statistics are extracted and formatted into a dictionary.
\end{itemize}

\item{\bf Real-Time Communication :}
\begin{itemize}
    \item \texttt{Flask-SocketIO} is used to emit the parsed data to the frontend in real time. This allows the frontend to display the latest statistics as they are gathered by the backend.
    \item The data is emitted as a WebSocket event (\texttt{'pidstat\_data'}), making it available for frontend visualization.
\end{itemize}
\end{itemize}
\subsection{Process Migration}
\subsubsection{Enabling tracking to log the process migration}
\begin{itemize}
    \item Navigate to the tracing directory. \mintinline{bash}{cd /sys/kernel/debug/tracing}
    \item Enable "sched\_migrate\_task" that allows to track and log the migration of tasks.
    
    \mintinline{bash}{echo 1 | sudo tee events/sched/sched_migrate_task/enable}
    \item Start tracing the events. \mintinline{bash}{echo 1 | sudo tee tracing_on}
    \item Wait for some time and let tracer log some migrations.
    \item Stop the tracing. \mintinline{bash}{echo 0 | sudo tee tracing_on}
    \item Check the trace file. \mintinline{bash}{cat trace}
\end{itemize}
\subsubsection{Logged Data}
\begin{itemize}
    \item An example of the logged data is shown below.
    \begin{minted}{text} 
        Chrome_ChildIOT-33594   [003] d..2. 27300.882725: sched_migrate_task: comm=Compositor pid=32759 prio=120 orig_cpu=0 dest_cpu=3
    \end{minted}
    \item 
\end{itemize}
\subsection{Studying the proc folder}
\subsubsection{What is the proc folder?}
\subsubsection{Add the folder name you studied!!}
\begin{itemize}
    \item try to put important points in bullet points
    \item At the end add some points on what you feel can be done with that information.
\end{itemize}
\end{document}